%% ====================================================================
%%  Chapter 2 — Literature Review
%% ====================================================================
\chapter{Literature Review}
\label{ch:literature}

This chapter covers the relevant literature on indoor localization and knowledge distillation, efficient
deep learning models, and robustness of deep learning models to signal degradation. Each of these topics
is covered in its own section, and each section ends with a discussion of the relevance of that topic to
this thesis.

\section{WiFi Indoor Localization}
\label{sec:lit:localization}

Indoor localization has been studied as a technology for the past two decades, mainly as a result of the
fact that GPS signals do not effectively reach individuals who are within indoor environments.
Technologies such as WiFi and Bluetooth have been investigated as potential replacements for GPS within
indoor environments.

WiFi localization is one of the most studied technologies for indoor positioning. It uses infrastructure
that is already installed in most buildings. Each WiFi access point (AP) continuously emits signals that
contain Channel State Information (CSI). CSI describes the strength of the signal on each of the
subcarriers continuously transmitted by each AP. The strength of the signal on each subcarrier is
related to the distance between the AP and the WiFi device; thus, by measuring the CSI of each AP, a
WiFi device can be localized to within sub-meter accuracy. This method was realized as a deep learning
model named DLoc that achieved sub-meter median accuracy~\cite{ayyalasomayajula2020deep}. In addition to
this model, the dataset (named WILD) upon which it was trained was also
published~\cite{ayyalasomayajula2020deep}.

\paragraph{Relevance to this work.} DLoc and the WILD dataset are used directly throughout this work:
DLoc is both the teacher and the accuracy baseline for our compressed models, and WILD is the dataset on
which every experiment is conducted.

\section{Knowledge Distillation}
\label{sec:lit:distillation}

Knowledge distillation (KD) was introduced by Hinton et al.~\cite{hinton2015distilling} as a technique
by which a small ``student'' neural network can be trained to learn the same behavior as a much larger
``teacher'' network. Knowledge-distillation techniques train the student not only to minimize the
difference between its predictions and the ground-truth labels of the training data, but also to
minimize the difference between its soft predictions and the soft labels of the teacher. The soft labels
of the teacher are typically generated by applying a softmax activation function to the teacher's output
logits, which transforms the teacher's raw outputs into a probability distribution over the possible
classifications. In classification problems, the predictions and labels of each training example are the
one-hot encoded label of the true class. The soft labels, in contrast, contain a probability that each
class is the true class of the example. These probabilities convey information that the teacher network
has learned about the similarities between the classes, so using them as the target of the student
enables the student to learn those similarities as well. Distillation has been applied widely to
classification problems, but it has also been applied to problems in which the target is not a
classification; the benefit that the soft labels provide in those settings differs from the benefit they
provide for classification.

\paragraph{Relevance to this work.} Knowledge distillation is the mechanism we use to transfer the
localization behavior of the full DLoc teacher into our lightweight students, so that a much smaller
student can learn to localize to within the accuracy of the teacher. As Chapter~\ref{ch:results} shows,
for our spatial-regression task this transfer turns out to act primarily on robustness rather than on
in-distribution accuracy.

\section{Efficient Neural Architectures}
\label{sec:lit:architectures}

The four models that comprise the student networks each incorporate ideas from different lines of
existing computer-vision and deep-learning research. The majority of these ideas were introduced as
efficient alternatives to conventional deep learning models, either in terms of their required
computational resources or their memory footprint. Each is described in the sections that follow.

\subsection{Residual Learning}
\label{sec:lit:resnet}

ResNets~\cite{he2016resnet} were among the first deep learning models to incorporate the idea of ``skip
connections'' into their architectures; these connections allowed gradients to flow through deep
networks to each of the network's parameters, enabling the training of deep learning models with many
hidden layers. The success and widespread adoption of ResNets demonstrated the value of such
``shortcuts'' in deep learning models, an idea that is also implemented in the skip connections used in
this thesis.

\subsection{MobileNetV2}
\label{sec:lit:mobilenet}

MobileNetV2~\cite{sandler2018mobilenetv2} was introduced as a mobile deep learning model that could be
deployed on devices with limited computational resources. MobileNet models replace conventional
convolution layers with depthwise-separable convolutions, which reduce the number of multiply--accumulate
operations required by roughly an order of magnitude; they also employ inverted residual blocks, which
enable MobileNet models to maintain a small memory footprint relative to conventional deep neural
networks. MobileNetV2 is used as the encoder for three of the four students proposed in this thesis.

\subsection{U-Net Skip Connections}
\label{sec:lit:unet}

U-Net models~\cite{ronneberger2015unet} were first applied to problems in the field of biomedical image
analysis, where deep learning models had to segment an image into different regions. The architecture of
a U-Net includes skip connections that enable the model to directly access high-resolution image
features that would otherwise be lost during the model's encoding process. U-Net models are therefore
naturally applicable to the problem of generating a heatmap from CSI measurements, where the system must
accurately determine the spatial location of the device emitting the WiFi signal. These skip connections
are incorporated into the MobileNetV2-UNet model in Chapter~\ref{ch:results}.

\subsection{State-Space Models (Mamba)}
\label{sec:lit:mamba}

State-space models have recently emerged as an alternative to models based on self-attention, both
because of their efficiency in terms of computational and memory requirements and because they have
exhibited strong performance on sequence-modeling tasks. Mamba models~\cite{gu2023mamba} employ a new
type of attention-like mechanism whose parameters are adapted based on the content of the input, and
have been applied to computer-vision tasks~\cite{visualmamba2024survey}. The inclusion of a student model
based on the Mamba architecture explores whether an inductive bias different from that of conventional
deep learning models is competitive against those alternatives.

\section{Robustness to Degraded and Shifted Distributions}
\label{sec:lit:robustness}

In addition to measuring the accuracy of a model on high-quality data, it is often important to evaluate
the model's accuracy in the face of degraded or differently distributed data; these measures of a
model's robustness are particularly important for models that are to be deployed in real-world
environments. For vision-based deep learning models, one way of measuring robustness is to introduce
degradations to the test images, such as noise, blur, weather effects, or digital image-processing
artifacts. Hendrycks and Dietterich~\cite{hendrycks2019robustness} published a standard deep learning
benchmark for evaluating robustness to these types of image corruption. Relatedly, it is often important
to determine how well a model generalizes to data that differs from the training data in ways that are
not explicitly accounted for during training, such as changes in the environment over time.

\paragraph{Relevance to this work.} The robustness of a model's predictions to such corruptions and
distribution shifts is particularly important for WiFi localization systems: APs can fail, there can be
interference in the WiFi signal, and the physical layout of the area being localized can change over
time. We therefore evaluate the robustness of each student to these corruptions and shifts, in the
spirit of~\cite{hendrycks2019robustness}, and we treat the WILD August sessions as a graded,
physically grounded distribution-shift benchmark.
